\documentclass[pre,twocolumn,superscriptaddress,nofootinbib]{revtex4-2}
\usepackage{graphicx}
\usepackage{amsmath}
\usepackage{booktabs}
\usepackage[hidelinks]{hyperref}

\begin{document}

\title{Curled tori with six and eight links: an additive ladder of directions, walls that rise at every rung, a
first move whose price is fixed with size, and a second direction that does not follow the first}

\author{Emily Smith}
\email{emilysmithcreate@gmail.com}
\affiliation{Independent researcher, Charlottesville, Virginia, USA}

\date{DRAFT, 25 September 2026}

\begin{abstract}
We extend the curled-torus study of a companion paper \cite{paper1} from four links per vertex to six and eight, in
the family of graph energies obtained from combinatorial quantum gravity (CQG) by raising the coefficient
$\lambda$ of its local term above 1 (only $\lambda = 1$ is CQG; everything here is at $\lambda > 1$, next to CQG and
not in it). Exact: the ladder of curled directions is additive at every number of links, $4(\lambda - 1)$ per vertex
per curled direction; the cheapest single move out of each rung, listed exhaustively (with eight links, over nearby
partners, a search checked against the full one at six), rises in cost from rung to rung (with six links 16, 36;
with eight, 20, 40, 80, at $\lambda = 1.25$), so nothing in the energy ties the directions together; and the cheapest
move out of flat space costs 32, 64 and 128 with four, six and eight links, the same at every $\lambda$. Measured,
pre-registered, sealed and energy-conserving: the price of the first move out of a six-link torus with two directions
curled is exactly 16 at every size from 192 to 512 vertices, no torus leaving its start below it and every one at or
above it; given that push, the first curled direction opens completely in 28 of 84 runs, most of the way in 21 and
stalls early in 32, and the second never opens in a cold bath, so the tori's verdict, read as registered, is NEVER
OPENS; a gas of fully curled 6-cubes at $\lambda = 1.10$, reported separately, joins and opens two of its three
directions in 21 of 24 replicas, without resting between them, but never into one space; and with eight links a push
equal to the first move's price buys that one move and nothing more. The author's rule that the space directions curl
and open together predicted a cascade; in a torus they do not open together. The published six-link curve has not
been reproduced, there is no published eight-link curve, and every result here carries that caveat.
\end{abstract}

\maketitle

\section{Introduction}

The companion paper \cite{paper1} studied a four-link torus with one direction curled, the two-dimensional model's
stand-in for a curled arrangement, and found a metastable state that opens into flat space by nucleation and a single
front. The author's framework holds that the real curled arrangement has three, or four, directions curled together,
that they open together, and that one activation releases the whole burp \cite{repo}. That rule cannot be posed with
four links, where a fully curled state (a gas of 4-cubes) has no wall at any $\lambda > 1$ and falls apart at once
\cite{repo}. With six and eight links it can, and this paper poses it: exactly, by listing the single moves out of
every rung of the ladder; and by sealed runs, each pre-registered with the author's prediction.

Two caveats govern everything below. The six-link kernel is the four-link code with the number of links read from
the array; at four links it makes the companion paper's chain draw for draw, its energy and moves are checked by
brute force against an independent implementation, and the ladder's energies are exact. But the published six-link
curve \cite{T22} has not been reproduced by it, under any reading of that figure's axis we have tested
(Sec.~\ref{sec:gate}); the author decided to proceed while that stays open, and no statement here is a statement about
the published model. With eight links there is no published curve at all, and the kernel is validated at four and six
links only. And the model has no time: a direction of a torus is a direction, and which one would be time
the model cannot say.

\section{Exact results}

\subsection{The ladder}

With $D$ directions, so $2D$ links per vertex, $S$ squares and $X = \sum_e (S_e - (2D - 2))_+$ surplus squares,
\begin{equation}
  H = 16\Big(\tfrac{D(D-1)}{2} N - S\Big) + 4\lambda X ,
\end{equation}
derived by us from the published definitions as the four-link energy was, as we read them \cite{repo}, and not yet
validated against a published curve. The flat $D$-torus (every
side at least 6) has $H = 0$. A side curled to length 4 puts one extra square on every edge along it, so each curled
direction costs $4(\lambda - 1)$ per vertex, additively: with six links the rungs are $0$, $4(\lambda - 1)$,
$8(\lambda - 1)$ and $12(\lambda - 1)$ (the 6-cube); with eight links there is a fourth rung at $16(\lambda - 1)$
(the 8-cube). The local dimension $d(v)$, the number of edge pairs at $v$ that close no square, counts the open
directions: $D$ on the flat torus, one less per curled direction, above $D$ where squares are missing.

\subsection{The walls}

Every switch the chain can propose out of each curled six-link torus was listed and priced at
$-16\,\Delta S + 4\lambda\,\Delta X$; by translation symmetry one vertex serves as the first of the pair. With eight
links, and for flat six-link space, only switches with nearby partners were listed, a shortcut checked against the
full search on one six-link torus ($4 \times 4 \times 18$) \cite{repo}. Table~\ref{tab:walls} gives the
cheapest kind for every rung at the sizes used. Three things follow. The walls rise from rung to rung: with six
links 16 then 36, with eight links 20, 40, 80 at $\lambda = 1.25$ (each lower wall is less by $16\lambda$ or
$32\lambda$, so the factor of two at eight links holds at this $\lambda$ only). So the energy ties nothing together; each direction has its own wall,
and the second is the higher. The window of $\lambda$ in which a rung is stuck for now (its cheapest single move out
costs energy, a property of the energy alone, unlike the chain-dependent ``stuck'' of the four-link map
\cite{paper1}) widens as more directions are open (with eight links, 1.25 for the gas, 1.43, 1.67, 2.5 up the
ladder). And the wall of flat space itself, the cheapest move out of the perfect lattice, is 32, 64 and 128 with four,
six and eight links, independent of $\lambda$; with eight links it sits above every rung's wall, as it does with four
(32 against the tube's 12). Whether a bath hot enough to pay the later rungs leaves the flat state intact is not
settled by these walls: melting is a matter of free energy, not of the cheapest move, flat three-dimensional space,
whose wall is 64, is melted by $g = 10$ in an exploratory scan \cite{repo}, and the ratio of flat space's wall to the
last rung's is smaller with eight links (128 against 80) than with four. One correction to the
project's earlier record: the six-link two-curled torus with a short open side of 6 admits a move at $96 - 80\lambda$
that a long open side does not, so its window is $1 < \lambda < 1.2$ at that size and $1 < \lambda < 1.5$ at the sizes
run here \cite{repo}.

\begin{table}[b]
\caption{The cheapest single move out of each curled state, exact. Cost as a function of $\lambda$ and at 1.25;
the window is where the cost is positive. Six-link tori $4 \times 4 \times L$ ($L \ge 18$), $4 \times L \times L'$;
eight-link tori at $N = 2304$.}
\label{tab:walls}
\begin{ruledtabular}
\begin{tabular}{llllc}
links & curled & $(\Delta S, \Delta X)$ & cost & at 1.25 \\
\hline
4 & one (the tube) & $(-2, -4)$ & $32 - 16\lambda$ & 12 \\
4 & flat sheet & $(-2, 0)$ & 32 & 32 \\
6 & one & $(-6, -12)$ & $96 - 48\lambda$ & 36 \\
6 & two & $(-6, -16)$ & $96 - 64\lambda$ & 16 \\
6 & three (6-cube, gas) & $(-6, -20)$ & $96 - 80\lambda$ & $-4$ \\
6 & flat 3-torus & $(-4, 0)$ & 64 & 64 \\
8 & one & $(-10, -16)$ & $160 - 64\lambda$ & 80 \\
8 & two & $(-10, -24)$ & $160 - 96\lambda$ & 40 \\
8 & three & $(-10, -28)$ & $160 - 112\lambda$ & 20 \\
8 & four (8-cube, gas) & $(-10, -32)$ & $160 - 128\lambda$ & 0 \\
8 & flat 4-torus & $(-8, 0)$ & 128 & 128 \\
\end{tabular}
\end{ruledtabular}
\end{table}

\section{Sealed runs with six links}

The sealed dynamics are the companion papers' \cite{paper1,paper2} on the six-link graphs: a bath of $C$ stores,
one holding the spark and the rest empty, a proposed switch drawing on one store at random and refused if it cannot
pay, $H$ plus the stores conserved to the last unit. The starting torus is $4 \times 4 \times L$: two directions
curled, one open, $8(\lambda - 1)$ per vertex above the flat 3-torus of the same $N$ ($6 \times 6 \times 8$ at
$N = 288$, $8 \times 8 \times 8$ at 512).

\subsection{The price of the first move is fixed with size}

\emph{Pre-registered as T32, the author predicting a fixed wall.} With a single store holding the spark, at
$N = 192$, 288, 384 and 512 ($L = 12$ to 32), sparks of 12, 14, 15, 16, 17, 18 and 20, ten replicas each, 20,000
sweeps: no replica left the start at 12, 14 or 15, and every replica left at 16 and above, at every size. A replica
counts as leaving when its squares change at all, so the threshold is the price of the first move, exactly the wall
of Table~\ref{tab:walls}, and it does not grow with the system. With four links the same price, 12, also sufficed to
convert the torus \cite{paper1}; with six links it does not always suffice (below), and with eight it did not
(Sec.~\ref{sec:eight}).

\subsection{Given the push, the second direction stays curled in a cold bath}

\emph{Pre-registered as T30, the author predicting that both curled directions open together (ALL AT ONCE); ours,
that the first opens and the second waits (FIRST ONLY).} At $\lambda = 1.25$, spark 16, twelve replicas per cell,
100,000 sweeps, seven cells ($C = 2N$, $N$, $N/2$, $N/4$ and $N/8$ at $N = 288$; $C = 2N$ and $N/4$ at 512), 84
replicas: the first curled direction opens completely, no vertex left at $d = 1$, in 28 of them; in 21 it goes most
of the way, 15 to 24\% of vertices still at $d = 1$; in 32 it stalls early, 36 to 53\% still at $d = 1$ and the
energy 1.86 to 1.95 per vertex against 2.0 at the start and 1.0 on the rung, so 5 to 14\% of the first rung's
release; the bath ends near 0.4 to 0.5 per store. The second direction never opens in a cold bath, its wall of 36
being unpayable there. With $C = N/8$ the bath heats to 1 to 5 per store and two of twelve replicas descend both rungs
to a defective near-flat state (88 to 90\% of vertices at $d = 3$, 3\% of squares lost), just under the
pre-registered cleanliness bar; the rest end as mixtures. At $\lambda = 1.10$ a spark of 26 (the wall is 25.6) buys
one move, which neither heals nor grows in 100,000 sweeps. No torus cell has a flat majority or a majority on the
one-curled rung, so, read as registered, the verdict for the tori is NEVER OPENS: the author's prediction failed, and
so did ours. The gas below is read by the same rules and its verdict is reported separately.

\subsection{The gas of 6-cubes}

The literal form of the author's rule is a state with no direction singled out: a gas of eight 6-cubes ($N = 512$),
stuck for now only below $\lambda = 1.2$, run at 1.10 with a spark of 8 (the wall) and 10, $C = 2N$, twelve replicas
each. In 21 of 24 replicas the cubes joined and opened two of their three directions, the energy falling from 1.2
to 0.24 to 0.41 per vertex, with about 40\% of vertices fully open in many small flat patches (the largest 36 to 64
vertices) and never one flat space; three replicas did not leave. Traced block by block, none rested at the
one-opened level for the pre-registered 5,000 sweeps (the longest stay was 3,800), though the whole descent took
10,000 to 46,000 sweeps. By the letter its verdict, T30-gas, is FIRST ONLY, which for the gas means two of three
directions open. This is further than the two-curled torus went at either $\lambda$: at 1.10 its first move led
nowhere, and at 1.25 its second wall (36) was not paid. Why the gas went further is not settled: its first wall is
lower (8), and the moves that join two cubes have not been priced.

\section{Eight links: the pattern test}
\label{sec:eight}

\emph{Pre-registered as T33.} With eight links the author's two variants can be told
apart: four directions tied (all open together), or a singleton beside three tied ones (one opens alone, then three
together); ours, one at a time. At $N = 2304$ every rung exists. Two starts, sealed, given exactly their wall: a
$4 \times 4 \times 4 \times 36$ torus at $\lambda = 1.25$ (spark 20) with $C = 2N$, $N/2$, $N/4$ and $N/8$, and a gas
of nine 8-cubes at $\lambda = 1.10$ (spark 20, the wall being 19.2) with $C = 2N$ and $N/4$; six replicas each,
50,000 sweeps. The pattern is read from the sequence of majority rungs (the $d$ held by more than half the vertices)
and the rests on them.

Every replica stalled. From both starts and at every bath size, the spark bought exactly one move: the energy rose
by that move's cost (from 3.000 to 3.009 per vertex from the torus, from 1.600 to 1.608 from the gas), the bath was
left at zero, and nothing changed again in 50{,}000 sweeps, so no move out of the state one move away is downhill.
The verdict by the letter is NO CASCADE from both starts, and neither the author's tied pattern nor our one at a time
was tested, because nothing opened. What the run shows is about the push: with four links the push equal to the
cheapest first move (12) is exactly the push that starts the change, and with six links at $\lambda = 1.25$ it opens
the first direction fully in a third of the runs, but with eight links it bought one move and nothing more, so the
change needs more than its first move. The contrast is not only in the number of directions: with six links at
$\lambda = 1.10$ the torus also stalled after one move. How high the pass is has not been measured; the cost of the
second move from the stalled states is being computed. A scan of the push from 20 to 160 (T40, pre-registered) is the
test that finds it and then reads the pattern.

\section{The reproduction gate}
\label{sec:gate}

Rule 2 of this project is to reproduce a published curve before interpreting new results built on the code. The
six-link target is Fig.~3 of Ref.~\cite{T22}: squares per vertex against $\ln \hbar g$ at $N = 500$, $\lambda = 1$.
Our kernel does not reproduce it. Under the plain reading of the axis ($\hbar g$ our $g$, which the paper's Eqs.~(1)
to (3) support), our curve at $N = 500$ is fully ordered at every coupling in the published range with long
equilibration, and drops abruptly with hysteresis at a larger coupling; under a division by $N^{1/3}$ it melts too
early; a literal reading of the paper's Eq.~(3) in which the weight goes as $1/g^2$, which halves the curve's extent
on the log axis, was pre-registered and failed, and the same code's two-dimensional figure is not compressed. Noticed
after that run and not scored: a plain factor of two in the coupling brings the ordered side and the top of the
transition close (on the two cooling legs the crossings of 9 and 6 squares per vertex fall within 0.01 to 0.07 and
0.03 to 0.04 of the published values, inside the 0.10 by which the pre-registered reading was scored); what remains
different is the hot tail, which the published curve reaches a factor of 2.8 sooner in coupling, and the plateau
(10.07 against our 11.4). Two readings are open. The published graphs may carry triangles and pentagons, as the paper
says its ground states can; against that, the published hot end (1.09 squares per vertex) sits at our bipartite value
(1.15), where such graphs would sit near half of it. Or the published runs were short and out of equilibrium: in an
exploratory run, a fast heating leg from a slowly cooled state comes near the published plateau and steepness in our
code (a width of 0.76 against 0.645, suggestive and not a match), though it then needs a coupling factor of about 5
to 6, not 2, to line up, so the two observations are not yet one account. The paper states no protocol; we think the
second reading the better supported, and it is unverified. Until it is settled, everything above is a statement about our family and its kernel, validated on its
own terms, and not about the published model.

\section{Discussion}

Three exact facts and two measured ones. The ladder is additive at every number of links; the walls rise from rung
to rung; flat space's own wall is 32, 64 and 128, the same at every $\lambda$. The price of the first move out is
fixed with size in three dimensions as in two. And given that push, a torus's directions do not open together: the
first opens fully in a third of the runs and stalls early in more than a third, and the second sits behind a wall the
first release does not pay in a cold bath; the gas of cubes, where no direction is singled out, opens two of three
without a rest between them. The author's rule that the
directions curl and open together therefore describes an ingredient this energy does not contain, which the map now
says plainly. The eight-link test stalled after one move from both starts and tested no pattern; whether a
singleton-then-three pattern, which she found natural on seeing the six-link result, appears from the fully curled
start is the question of the push scan (T40). What would carry the question further is a term that
couples the directions, which would be a new knob and needs its own decision before any run.

\begin{acknowledgments}
Code, analysis and text were drafted with Claude (Anthropic) in conversation with the author, who directed and
checked the work; no physicist has reviewed it. The repository holds every configuration, seed, result and
pre-registration: \url{https://github.com/EmilySmithCreate/RecreationalPhysics}.
\end{acknowledgments}

\begin{thebibliography}{9}
\bibitem{paper1} E. Smith, The curled torus burps: front-driven decompactification between ordered states in a
graph model of emergent geometry, draft (2026), \texttt{docs/papers/curled\_torus/}.
\bibitem{paper2} E. Smith, What the burp leaves behind, draft (2026), \texttt{docs/papers/relic/}.
\bibitem{T22} C. A. Trugenberger, Emergent time, cosmological constant and boundary dimension at infinity in
combinatorial quantum gravity, JHEP 04, 019 (2022); arXiv:2112.03778.
\bibitem{repo} \url{https://github.com/EmilySmithCreate/RecreationalPhysics}: VISION Updates 22 to 26;
PREREGISTRATION T30, T32, T33, T40, Gate C$'$; ASSUMPTIONS Q21, O41, O43, O45, O46, O49, O50, O51, O53, O54, O55, O62,
O64.
\end{thebibliography}

\end{document}
